
\sourcenote{Genre: core rule (tier 1) --- an ordinary mathematical proof at
Putnam-exam-solution detail. Every hypothesis of the rule is used below and
remarked on; the formal-anchor table maps each step to the Lean declaration
that carries it. This section also fixes the semantic setting that the
other equality and conversion sections cite.}

\textbf{Setting.}\enspace Fix a hypernet $H$. A \emph{valuation} $\sigma$
assigns each wire a tensor of its declared shape over its declared
dtype (an extended dtype admits infinite entries); the rational term
evaluation
$\mathrm{ev}_\sigma(t)$ is \emph{partial} (a wire holding an infinite
value of an extended dtype, or a broadcast-incompatible operation, yields
no value). A comparison atom $t_1 \bowtie t_2$ \emph{holds} at $\sigma$
when both sides evaluate and the evaluated tensors stand in the relation
$\bowtie$ \emph{elementwise under broadcast}: equal shapes compare entry
by entry; a rank-0 (scalar) tensor against a tensor of any shape compares
against every entry; two unequal positive-rank shapes never compare. A
sequent $\Phi \vdash \Psi$ is \emph{semantically valid} when every
admissible $\sigma$ satisfying $\Phi$ satisfies $\Psi$. Note two
consequences that drive everything below: equality atoms are
\emph{evaluation-level} facts, not syntactic ones; and a holding equality
between tensors of \emph{different} broadcast shapes (a scalar against a
vector) does not make the two sides interchangeable --- it only says the
scalar matches every entry.

\textbf{Claim} (\texttt{sound\_eqElim}).\enspace The rule's payload
supplies a hole wire $a$, terms $t, u$, a body formula $\Psi$, and checked
side conditions: the two premises share the composite premise $\Phi$; the
first premise concludes the atom $t = u$; the second premise concludes
$\Psi\sub{a}{t}$; the conclusion is $\Psi\sub{a}{u}$; and (the typing
condition) $\exists s.\ \bshape_H(t) = s = \bshape_H(u)$ --- both terms
carry a \emph{common defined} broadcast shape in $H$. If both premises are
semantically valid, so is $\Phi \vdash \Psi\sub{a}{u}$.

\begin{blemma}[substitution congruence, term level]\label{lem:eqelim-term}
Fix $\sigma$, and suppose $\mathrm{ev}_\sigma(t) = \mathrm{ev}_\sigma(u)$
as partial values (both undefined, or both defined and equal). Then for
every term $e$,
$\mathrm{ev}_\sigma(e\sub{a}{t}) = \mathrm{ev}_\sigma(e\sub{a}{u})$.
\end{blemma}
\begin{proof}
Structural induction on $e$. The base case $e = a$ is the hypothesis; the
base cases $e$ a constant or a different wire are reflexivity; every
compound constructor is a congruence of the evaluator (the evaluator of a
sum, product, difference, negation, or projection is a function of the
evaluations of the parts). Two remarks. First, the hypothesis and
conclusion are equalities of \emph{partial} values --- the lemma never
manufactures definedness, it only transports whatever is there; in the
theorem, definedness will come out of the equality premise. Second, the
term grammar has no binders, so substitution is plain leaf replacement
and no capture condition is needed anywhere.
\end{proof}

\begin{blemma}[substitution congruence, formula level]\label{lem:eqelim-form}
Under the hypothesis of Lemma~\ref{lem:eqelim-term}: for every formula
$F$,\ $\sigma \models F\sub{a}{t}$ iff $\sigma \models F\sub{a}{u}$.
\end{blemma}
\begin{proof}
Induction on $F$. Comparison atoms substitute in both sides and satisfy
the same relation by Lemma~\ref{lem:eqelim-term}; a region atom
substitutes in its term slot only (its simplex and extension-set data
carry no wires), so its holding --- an assertion about the evaluated
tensor --- transports the same way; $\wedge$ and $\vee$ are the inductive
steps; $\top$ and $\bot$ are untouched by substitution.
\end{proof}

\begin{blemma}[the shape pin]\label{lem:eqelim-pin}
Suppose $\sigma \models t = u$, and suppose $t$ and $u$ carry a common
defined broadcast shape $s$ in $H$. Then
$\mathrm{ev}_\sigma(t) = \mathrm{ev}_\sigma(u)$ --- equality of
\emph{tensors}, not merely the elementwise relation.
\end{blemma}
\begin{proof}
The holding atom yields evaluations $A$ of $t$ and $B$ of $u$ in the
elementwise-equality relation. Broadcast-shape soundness
(\texttt{Term.evalQ\_bshape}) pins $A.\mathrm{shape} = s =
B.\mathrm{shape}$: a defined evaluation of a $\bshape$-typed term lands
exactly at the declared broadcast shape. Elementwise equality at
\emph{equal} shapes is equality of the (equal-length) entry lists,
position by position, hence equality of the tensors. This is the one step
that consumes the typing condition, and it is false without it: the
elementwise relation between a rank-0 tensor and a vector says only that
the scalar matches every entry (Remark~\ref{rem:eqelim-cex}).
\end{proof}

\begin{bthm}[soundness of \rulename{EqElim}]\label{thm:eqelim-sound}
Under the payload's checked hypotheses, if $\Phi \vdash t = u$ and
$\Phi \vdash \Psi\sub{a}{t}$ are semantically valid, then so is
$\Phi \vdash \Psi\sub{a}{u}$.
\end{bthm}
\begin{proof}
Fix an admissible $\sigma$ with $\sigma \models \Phi$. By validity of the
first premise, $\sigma \models t = u$; with the typing condition,
Lemma~\ref{lem:eqelim-pin} upgrades this to
$\mathrm{ev}_\sigma(t) = \mathrm{ev}_\sigma(u)$ (both defined and equal,
so in particular equal as partial values). By validity of the second
premise --- its premise formula is $\Phi$ itself, by the checked premise
match --- $\sigma \models \Psi\sub{a}{t}$. Lemma~\ref{lem:eqelim-form}
transports this to $\sigma \models \Psi\sub{a}{u}$, which is the
conclusion formula by the checked substitution equation.
\end{proof}

\begin{bremark}[why the typing condition is not decoration]
\label{rem:eqelim-cex}
Drop the common-shape condition and the rule is \emph{unsound}, by a
counterexample recorded in the soundness development's header inventory
(orientation of the equality atom flipped here to spare an
\rulename{EqSym} step). Let $v$ be a wire of shape $[1]$ over $\Q$, let
$\Delta$ be the segment $\mathrm{conv}(1,2) \subset \Q^1$, and take
\[
\Phi \;:=\; \bigl(v = \tfrac{3}{2}\bigr) \;\wedge\; \bigl(v \in
\mathcal{R}(\Delta, \emptyset)\bigr).
\]
$\Phi$ is satisfiable: set $v := [\,3/2\,]$; the equality holds by rank-0
broadcast (the scalar $3/2$ matches the single entry), and the region
atom holds. \rulename{Ax} and $\wedge$E derive $\Phi \vdash v = 3/2$ and
$\Phi \vdash v \in \mathcal{R}(\Delta, \emptyset)$; the un-gated rule
with body $\Psi := (a \in \mathcal{R}(\Delta, \emptyset))$, $t := v$,
$u := 3/2$ then concludes
$\Phi \vdash (\tfrac{3}{2} : \mathrm{Term}) \in \mathcal{R}(\Delta,
\emptyset)$ --- which is \emph{unsatisfiable}: a constant term evaluates
to a rank-0 tensor, and a region atom demands a rank-1 tensor of extent
equal to the ambient dimension. A satisfiable premise, an unsatisfiable
conclusion. The root cause is precisely the gap between ``elementwise
equal under broadcast'' and ``interchangeable'': $3/2 = [\,3/2\,]$ holds
without the terms being exchangeable in region contexts. With the typing
condition the payload is uninhabitable here ---
$\bshape_H(\tfrac{3}{2}) = [\,]$ while $\bshape_H(v) = [1]$, and no
common $s$ exists --- so the certificate refuses. Note what the repair
does \emph{not} do: broadcast equalities remain derivable and usable as
facts; what is forbidden is exactly using one as a Leibniz license.
\end{bremark}

\begin{bremark}[the equational kit; relation to practice]
\label{rem:eqelim-kit}
\rulename{EqElim} with \rulename{EqRefl} and \rulename{EqSym} (next two
sections) completes a certificate-side equational kit: reflexivity and
symmetry are primitive, and transitivity is a one-node derivation ---
from $\Phi \vdash t = u$ and $\Phi \vdash u = w$, apply \rulename{EqElim}
with body $\Psi := (t = a)$, rewriting the $u$ in the \emph{first}
premise (taken as body premise, $\Psi\sub{a}{u}$) to $w$; the derived
node inherits the common-shape typing obligation on $(u, w)$ ---
broadcast equality is not transitive without it. On the producer side,
equational goals are typically discharged by congruence-closure engines
\cite{nelsonoppen1980,nieuwenhuisoliveras2007}; what ships in a
certificate is not the closure computation but its residue: a chain of
\rulename{EqRefl}/\rulename{EqSym}/\rulename{EqElim} nodes, each checked
by syntactic recomputation of one substitution.
\end{bremark}

\begin{anchors}
Setting: elementwise-under-broadcast atom semantics & \texttt{QTensor.rel2}, \texttt{Formula.holds} (Layer 10)\\
Lemma \ref{lem:eqelim-term} & \texttt{Term.evalQ\_subst\_congr}\\
Lemma \ref{lem:eqelim-form} & \texttt{Formula.holds\_subst\_congr}\\
Lemma \ref{lem:eqelim-pin} & \texttt{Term.evalQ\_bshape}; \texttt{QTensor.eq\_of\_rel2}; payload field \texttt{tyEq}\\
substitution equations & \texttt{substituteInFormula} = \texttt{Formula.subst} (payload fields \texttt{bodySubT}, \texttt{conclSubU})\\
premise sharing & payload fields \texttt{preEq}, \texttt{preMatch}\\
Theorem \ref{thm:eqelim-sound} & \texttt{sound\_eqElim}; wired as core \texttt{RuleData.eqElim} into \texttt{ProofObject.sound}\\
Remark \ref{rem:eqelim-cex} & the counterexample inventory, \texttt{Soundness.lean} header\\
\end{anchors}
